\documentclass[a4paper,12pt]{article}

%Acrônimos de periódicos da American Astronomical Society (AAS)
\newcommand\aj{Astronomical Journal}%        % Astronomical Journal 
\newcommand\psj{Planetary Science Journal}%       % Planetary Science Journal
\newcommand\araa{Annual Review of Astron and Astrophys}%  % Annual Review of Astron and Astrophys 
\newcommand\apj{Astrophysical Journal}%    % Astrophysical Journal 
\newcommand\apjl{Astrophysical Journal, Letters}%     % Astrophysical Journal, Letters 
\newcommand\apjs{Astrophysical Journal, Supplement}%    % Astrophysical Journal, Supplement 
\newcommand\ao{Applied Optics}%   % Applied Optics 
\newcommand\apss{Astrophysics and Space Science}%  % Astrophysics and Space Science 
\newcommand\aap{Astronomy and Astrophysics}%     % Astronomy and Astrophysics 
\newcommand\aapr{Astronomy and Astrophysics Reviews}%  % Astronomy and Astrophysics Reviews 
\newcommand\aaps{Astronomy and Astrophysics, Supplement}%    % Astronomy and Astrophysics, Supplement 
\newcommand\azh{Astronomicheskii Zhurnal}%       % Astronomicheskii Zhurnal 
\newcommand\baas{Bulletin of the American Astronomical Society}%     % Bulletin of the AAS 
\newcommand\icarus{Icarus}% % Icarus
\newcommand\jaavso{The Journal of the American Association of Variable Star Observers}%  % The Journal of the American Association of Variable Star Observers
\newcommand\jrasc{Journal of the Royal Astronomical Society of Canada}%   % Journal of the RAS of Canada 
\newcommand\memras{Memoirs of the Royal Astronomical Society}%  % Memoirs of the RAS 
\newcommand\mnras{Monthly Notices of the Royal Astronomical Society}%   % Monthly Notices of the RAS 
\newcommand\pra{Physical Review A: General Physics}% % Physical Review A: General Physics 
\newcommand\prb{Physical Review B: Solid State}% % Physical Review B: Solid State 
\newcommand\prc{Physical Review C}% % Physical Review C 
\newcommand\prd{Physical Review D}% % Physical Review D 
\newcommand\pre{Physical Review E}% % Physical Review E 
\newcommand\prl{Physical Review Letters}% % Physical Review Letters 
\newcommand\pasp{Publications of the American Society of the Pacific}%     % Publications of the ASP 
\newcommand\pasj{Publications of the Astronomical Society of Japan}%     % Publications of the ASJ 
\newcommand\qjras{Quarterly Journal of the Royal Astronomical Society}%   % Quarterly Journal of the RAS 
\newcommand\skytel{Sky and Telescope}%   % Sky and Telescope 
\newcommand\solphys{Solar Physics}% % Solar Physics 
\newcommand\sovast{Soviet Astronomy}% % Soviet Astronomy 
\newcommand\ssr{Space Science Reviews}% % Space Science Reviews 
\newcommand\zap{Zeitschrift fuer Astrophysik}%       % Zeitschrift fuer Astrophysik 
\newcommand\nat{Nature}%  % Nature 
\newcommand\iaucirc{IAU Circulars}% % IAU Cirulars 
\newcommand\aplett{Astrophysics Letters}%  % Astrophysics Letters 
\newcommand\apspr{Astrophysics Space Physics Research}% % Astrophysics Space Physics Research 
\newcommand\bain{Bulletin Astronomical Institute of the Netherlands}% % Bulletin Astronomical Institute of the Netherlands 
\newcommand\fcp{Fundamental Cosmic Physics}%   % Fundamental Cosmic Physics 
\newcommand\gca{Geochimica Cosmochimica Acta}% % Geochimica Cosmochimica Acta 
\newcommand\grl{Geophysics Research Letters}%  % Geophysics Research Letters 
\newcommand\jcp{Journal of Chemical Physics}%     % Journal of Chemical Physics 
\newcommand\jgr{Journal of Geophysics Research}%     % Journal of Geophysics Research 
\newcommand\jqsrt{Journal of Quantitative Spectroscopy and Radiative Transfer}%   % Journal of Quantitiative Spectroscopy and Radiative Trasfer 
\newcommand\nphysa{Nuclear Physics A}%     % Nuclear Physics A 
\newcommand\physrep{Physics Reports}%       % Physics Reports 
\newcommand\physscr{Physica Scripta}%        % Physica Scripta 
\newcommand\planss{Planetary Space Science}%  % Planetary Space Science 
\newcommand\procspie{Proceedings of the SPIE}%      % Proceedings of the SPIE 

\newcommand\actaa{Acta Astronomica}%  % Acta Astronomica
\newcommand\caa{Chinese Astronomy and Astrophysics}%  % Chinese Astronomy and Astrophysics
\newcommand\cjaa{Chinese Journal of Astronomy and Astrophysics}%  % Chinese Journal of Astronomy and Astrophysics
\newcommand\jcap{Journal of Cosmology and Astroparticle Physics}%  % Journal of Cosmology and Astroparticle Physics
\newcommand\na{New Astronomy}%  % New Astronomy
\newcommand\nar{New Astronomy Review}%  % New Astronomy Review
\newcommand\pasa{Publications of the Astron. Soc. of Australia}%  % Publications of the Astron. Soc. of Australia
\newcommand\rmxaa{Revista Mexicana de Astronomia y Astrofisica}%  % Revista Mexicana de Astronomia y Astrofisica

%% added feb 9, 2016
\newcommand\maps{Meteoritics and Planetary Science}% Meteoritics and Planetary Science
\newcommand\aas{American Astronomical Society Meeting Abstracts}% American Astronomical Society Meeting Abstracts
\newcommand\dps{American Astronomical Society/Division for Planetary Sciences Meeting Abstracts}% American Astronomical Society/Division for Planetary Sciences Meeting Abstracts


%Informações (Título, autor, resumo, abstract, palavras-chave, etc)
%%% EDITÁVEL%%%
%numero de autores máximo 3
\newcounter{numerodeautores}\addtocounter{numerodeautores}{4}% numero 1, 2 ou 3 
\newcommand{\titulonotatecnica}{Título da Nota Técnica}
\newcommand{\autornotatecnica}{Autor1}
\newcommand{\filiacaoautordanota}{Laboratório Nacional de Astrofísica}
\newcommand{\emailautor}{autor1@lna.br} %necessario apenas para o primeiro autor
\newcommand{\numeronotatecnica}{LNA/NT/202X-z}
\newcommand{\datanotatecnica}{Mês/ano}
\newcommand{\segundoautor}{Autor2}
\newcommand{\terceiroautor}{Autor3}
\newcommand{\quartoautor}{Autor4}

%Revisores em ordem alfabética
\newcommand{\revisorum}{Revisor 1}
\newcommand{\revisordois}{Revisor 2}
\newcommand{\revisortres}{Revisor 3 -- nomes dos revisores em ordem alfabética}

% Resumo %
\newcommand{\resumonotatecnica}{Aqui vai um resumo do trabalho a ser apresentado ao longo da nota técnica}

% Abstract %
\newcommand{\abstractnotatecnica}{\textit{Include here the abstract of your work}}

% Palavras-chave/keywords
\newcommand{\keywordsnotatecnica}{coloque aqui as palavras-chave para futura indexação. Exemplo: espectrógrafo; rede echelle; projeto óptico / \textit{spectrograph; echelle grating; optical design}}

%Configurações (pacotes, comandos)
\usepackage[utf8]{inputenc}
\ifnum\value{linguagemdocumento}=1
  \usepackage[english,brazil]{babel}
\else \ifnum\value{linguagemdocumento}=2
    \usepackage[brazil,english]{babel}
  \fi \fi
\usepackage[T1]{fontenc}  %%%Codificação de fontes - ç,ã,é ... etc.
%%%Fontes 
%\usepackage{lmodern} %%% Fonte padrão Computer Modern do LaTeX
\usepackage[sfdefault]{FiraSans} %% Fira Sans da marca do LNA
\renewcommand{\familydefault}{\sfdefault} %% Fira Sans da marca do LNA

\usepackage{amssymb,amsmath}
\usepackage{indentfirst}  %Indenta o primeiro paragráfo.
\usepackage{graphicx}
\graphicspath{{Fig/}}
\usepackage[final]{pdfpages}
\usepackage[final]{hyperref}
\usepackage{color}
\usepackage[normalem]{ulem}
\usepackage{longtable}
\usepackage{tabularx}
\setlength{\parindent}{1.25cm}
\usepackage{setspace}
\usepackage{array}
\usepackage{booktabs}
\usepackage{csquotes}

%Carregando o pacote caption
\usepackage[
  format=hang,
  font=small,
  labelfont=bf
]{caption}
%Formantando o caption de figuras
\captionsetup[figure]{%
  width=.8\linewidth%
}
%Formantando o caption de tabelas
%\captionsetup[table]{
%	position=above,
%	justification=raggedright,
%	width=.8\linewidth
%labelsep=newline, % <<< label and text on different lines
%singlelinecheck=false % <<< raggadright also when the caption is shorter
% than a single line
%}
\usepackage{subcaption}


%Mudança da contagem
\counterwithin{table}{section}
\counterwithin{figure}{section}

%Inclui uma subseção a mais e na sumário
\setcounter{secnumdepth}{5}
\setcounter{tocdepth}{5}

\usepackage[a4paper,left=3cm,right=2.5cm,top=2cm,bottom=2.5cm,head=1.5cm,headsep=.5cm,includeheadfoot]{geometry}

\hypersetup{colorlinks,
  debug=false,
  linkcolor=blue,  %%% cor do tableofcontents, \ref, \footnote, etc
  citecolor=blue,  %%% cor do \cite
  urlcolor=magenta,   %%% cor do \url e \href
  bookmarksopen=true,
  pdftitle={},
  pdfauthor={},
  pdfsubject={},
  pdfkeywords={}
}

%fancy package following by fancy commands declarations
\usepackage{fancyhdr}

\fancypagestyle{capa}{%
  \fancyhead{}%
  \fancyfoot{}%
  \lhead{}%
  \chead{}%
  \rhead{\includegraphics[scale=.22]{marca_LNA_horiz_grad_nome_pos_RGB.pdf}}%
  \lfoot{}%
  \cfoot{}%
  \rfoot{}%
  \renewcommand{\headrulewidth}{0pt} % remove lines as well
  \renewcommand{\footrulewidth}{0pt} % remove lines as well
}

\fancypagestyle{contracapa}{%
  \fancyhead{}%
  \fancyfoot{}%
  \lhead{}%
  \chead{}%
  \rhead{}%
  \lfoot{}%
  \cfoot{}%
  \rfoot{}%
  \renewcommand{\headrulewidth}{0pt} % remove lines as well
  \renewcommand{\footrulewidth}{0pt} % remove lines as well
}

\fancypagestyle{normal}{%
  \fancyhead{}%
  \fancyfoot{}%
  \lhead{}%
  \chead{}%
  \rhead{\includegraphics[scale=.22]{marca_LNA_horiz_grad_nome_pos_RGB.pdf}}%
  \lfoot{\titulonotatecnica}%
  \cfoot{}%
  \rfoot{\small\thepage}%
  \renewcommand{\headrulewidth}{1pt} % remove lines as well
  \renewcommand{\footrulewidth}{1pt} % remove lines as well
}

\usepackage{enumitem}
\setlist[description]{%
  topsep=30pt,               % space before start / after end of list
  itemsep=5pt,               % space between items
  %font={\bfseries\sffamily}, % set the label font
  %  font={\bfseries\sffamily\color{red}}, % if colour is needed
}

\usepackage[
  style=abnt,
  backref=true,
  backend=biber,
  citecounter=true,
  backrefstyle=three,
  url=true,
  maxbibnames=99,
  mincitenames=1,
  maxcitenames=3,
  backref=true,
  hyperref=true,
  giveninits=true,
  uniquename=false,
  scbib=false,
  indent=true,
  justify=true,
  noslsn,
  repeatfields,
  uniquelist=false]{biblatex}
%\let\UpperOrSC\NormalOrSC
\let\UpperOrSCCite\NormalOrSCCite%
\usepackage{xpatch}
\xpatchbibmacro{name:andothers}{%
  \bibstring{andothers}%
}{%
  \bibstring[\emph]{andothers}%
}{}{}
\addbibresource{XX/bibliografia.bib}

\begin{document}

%%% Capa %%%
%NÃO EDITÁVEL
%Primeira página
\thispagestyle{capa}
\hfill
\vspace{2cm}


\begin{center}
    {\Large Notas Técnicas do Laboratório Nacional de Astrofísica}

    \vspace{6cm}

    %Título da nota técnica - Editar no arquivo info.tex
    {\LARGE \textbf{\titulonotatecnica}}

    \hfill

    \vspace{1cm}

    %Autor da nota técnica - Editar no arquivo info.tex
    {\large \ifnum\value{numerodeautores}=1 \autornotatecnica%
        \else \ifnum\value{numerodeautores}=2 \autornotatecnica, \& \segundoautor%
            \else \ifnum\value{numerodeautores}=3 \autornotatecnica, \segundoautor, \& \terceiroautor%
                \else \ifnum\value{numerodeautores}=4 \autornotatecnica, \segundoautor, \terceiroautor, \& \quartoautor%
                    \fi \fi \fi \fi
    }

    \vfill

    %Numero da nota técnica - Editar no arquivo info.tex
    \Large \numeronotatecnica%

    %Data da nota técnica - Editar no arquivo info.tex
    \Large \datanotatecnica%

\end{center}


\pagestyle{normal}
%%% Contracapa %%%
\newpage

\begin{center}

    \hfill

    \vspace{1cm}

    {\Large \textbf{\titulonotatecnica}}

    \hfill
    \vspace{.5cm}

    {\large \ifnum\value{numerodeautores}=1 \autornotatecnica\textsuperscript{1} %email apenas do primeiro autor 
        \else \ifnum\value{numerodeautores}=2 \autornotatecnica\textsuperscript{1}, \&  \segundoautor%
            \else \ifnum\value{numerodeautores}=3 \autornotatecnica\textsuperscript{1}, \segundoautor, \& \terceiroautor%
                \else \ifnum\value{numerodeautores}=4 \autornotatecnica\textsuperscript{1}, \segundoautor, \terceiroautor, \& \quartoautor%
                    \fi
                \fi
            \fi
        \fi
    }

    \hfill
    \vspace{.5cm}

    \(^1\)\textit{\filiacaoautordanota, \emailautor}

    \vspace{1.5cm}

    \begin{minipage}{12cm}
        \textbf{Resumo:} \resumonotatecnica.%
    \end{minipage}

    \vspace{1.5cm}

    \begin{minipage}{12cm}
        \textbf{Abstract:} \abstractnotatecnica.%
    \end{minipage}


    \vspace{1.5cm}

    \begin{minipage}{12cm}
        \textbf{Palavras-chave/\textit{keywords}:} \keywordsnotatecnica.%
    \end{minipage}

    \vspace{1.5cm}

    \begin{minipage}{12cm}
        \textbf{Submetido em:}  \datanotatecnica.%
    \end{minipage}

    \vspace{1cm}

    \begin{minipage}{12cm}
        \textbf{Revisado por:} \\
        {\raggedleft%
            \renewcommand{\arraystretch}{1.0}%
            \begin{tabular}{ll}
                \hspace{.8cm} &                \\
                \hspace{.8cm} & {\revisorum}   \\
                \hspace{.8cm} &                \\
                \hspace{.8cm} & {\revisordois} \\
                \hspace{.8cm} &                \\
                \hspace{.8cm} & {\revisortres} \\
            \end{tabular}}
    \end{minipage}


\end{center}

\newpage


%%% Tabela de conteúdos %%%
\tableofcontents
%\listoffigures
%\listoftables

%\newpage
%Inclusão de acrônimos

% Acronimos e definicoes
\section*{Acrônimos e definições}
\addcontentsline{toc}{section}{Acrônimos e definições}

\noindent
{\renewcommand{\arraystretch}{1.25}% Espaçamento vertical entre as linhas
\begin{tabularx}{0.96\linewidth}{@{}>{\raggedright\arraybackslash}p{0.12\linewidth} >{\raggedright\arraybackslash}X@{}}
\textbf{LNA} & Laboratório Nacional de Astrofísica \\
\textbf{OPD} & \textsl{Observatório Pico dos Dias} \\
\end{tabularx}}


%Escopo do documento
\onehalfspacing%
%% Escopo do documento
\section*{Escopo do documento}
\addcontentsline{toc}{section}{Escopo do documento}

Escrever o escopo do documento.



%Texto da nota técnica
%\newpage

\section{Introdução}\label{sec:introducao}
\onehalfspacing%

Introdução do trabalho \(\ldots\)%

\section{Inclua quantas seções julgar necessárias}\label{sec:section1}

Desenvolvimento Tabela~\ref{tab:obs_nights} da seção~\(\ldots\)~\cite{mesquita2016}.


\begin{table}[ht!]
    \centering
    \caption{Example of a table.}%
    \label{tab:nights}
    \begin{tabular}{cc}
        \toprule
        Obs date     & N images \\
        \midrule
        2026--07--02 & 9243     \\
        2026--07--14 & 4857     \\
        2026--07--15 & 3992     \\
        2026--07--19 & 7711     \\
        2026--07--31 & 10546    \\
        2026--08--03 & 3821     \\
        2026--08--05 & 6611     \\
        2026--08--19 & 4304     \\
        2026--08--22 & 8973     \\
        2026--08--23 & 4014     \\
        \bottomrule
    \end{tabular}
\end{table}

\subsection{Inclua quantas subseções julgar necessárias}\label{subsec:subsection1}

\begin{table}[h!]
    \centering
    \caption{Example of a table.}%
    \label{tab:obs_nights}
    \begin{tabular}{cc}
        \toprule
        Obs date     & N images \\
        \midrule
        2026--07--02 & 9243     \\
        2026--07--14 & 4857     \\
        2026--07--15 & 3992     \\
        2026--07--19 & 7711     \\
        2026--07--31 & 10546    \\
        2026--08--03 & 3821     \\
        2026--08--05 & 6611     \\
        2026--08--19 & 4304     \\
        2026--08--22 & 8973     \\
        2026--08--23 & 4014     \\
        \bottomrule
    \end{tabular}
\end{table}

Desenvolvimento da subseção \(\ldots\)~\cite{einstein1905}

\section{Conclusões}\label{sec:conclusoes}

Sumarize os principais resultados e conclusões do desenvolvimento~\cite{latexcompanion}.


\section{Agradecimentos}

Inclua os agradecimentos que julgar necessários.

\vspace{0.5cm}

Não se esqueça de agradecer as agências de fomento que apoiaram o projeto e/ou autor (es), caso hajam!





% Para criar entrada no sumário sem inserir página/folha extra, use \phantomsection + \addcontentsline
%\phantomsection
%\addcontentsline{toc}{section}{Referências}
% Referências bibliográficas - obrigatório
\renewcommand{\baselinestretch}{1}
\setlength\bibitemsep{12pt}
\ifnum\value{linguagemdocumento}=1
    \printbibliography[heading=bibintoc,title={Referências}]
\else \ifnum\value{linguagemdocumento}=2
        \printbibliography[heading=bibintoc,title={References}]
    \fi \fi

%Bibliografia da nota técnica
%\newpage
%\begin{thebibliography}{9}
\addcontentsline{toc}{section}{Referências}

%Escrever conforme modelos abaixo
%Usar o comando \cite no texto para chamar a citação

\bibitem{SIFS}
  Laboratório Nacional de Astrofísica,
  \emph{SIFS},
  \url{http://www.lna.br/~sifs/index.html},
  acessado em: novembro de 2016.
  
\bibitem{PARRY}
	Parry, I.R.,
	\emph{The Astronomical Uses of Optical Fiber},
	Fiber optics in Astronomy III, 
	ASP Conference Series, Vol. 152,
	1998.

\bibitem{PFS}
	Institute for the Physics and Mathematics of the Universe,
  	\emph{Prime Focus Spectrograph - PFS},
  	\url{http://pfs.ipmu.jp/index.html},
  acessado em: agosto de 2019.
  
\bibitem{ENGIS}
	Engis Corporation,
	\url{http://www.engis.com/},
  acessado em: novembro de 2019.
  
\bibitem{MACANHAN}
	 Macanhan, V. B. P.; Arruda, M. V.,
	 Patente: Privilégio de Inovação. 
	 Número do registro: BR1020180754114, 
	 título: ``Atuador de precisão com mola de compressão'', 	
	 Instituição de registro: INPI - Instituto Nacional da Propriedade Industrial. 
	 Depósito: 07/12/2018
  

\end{thebibliography}


%Apêndice se necessário
\appendix
%\section{Título do Apêndice}

%Exemplo de inclusão de figura
%\begin{figure}[htb!]
%	\centering
%	\includegraphics[width=1\linewidth]{../XX/Fig/nome_da_figura.jpg}
%	\caption{Legenda da figura.}
%	\label{fig:nome_da_figura}
%\end{figure}

\end{document}